% Standalone manuscript insert.
% Requires \usepackage[breakable,skins]{tcolorbox} in the manuscript preamble.
% The `breakable` option lets the frame split cleanly at page boundaries.
\clearpage
\begin{tcolorbox}[
  breakable,
  enhanced,
  colback=white,
  colframe=black,
  boxrule=0.7pt,
  arc=0pt,
  left=10pt,
  right=10pt,
  top=8pt,
  bottom=8pt,
  before skip=12pt,
  after skip=12pt
]

\textbf{Box 3 \,|\, Extending location--scale models:\newline
Where does heterogeneity occur?}

The location--scale model in
Equations~\eqref{eq:locscale_loc}--\eqref{eq:locscale_scale} allows residual
heterogeneity among effect sizes within studies to depend on moderators. In
multilevel meta-analysis, however, heterogeneity can arise at more than one
level. Studies can differ in their average effects (between-study
heterogeneity), while effect sizes can also vary around those study-level
averages (within-study heterogeneity). Extending the location--scale framework
to both levels therefore allows us to ask not only whether heterogeneity changes
with context, but also \emph{where that heterogeneity occurs}.

Using the notation introduced above, a general multilevel location--scale model
can be written as
\begin{align}
  y_i
  &=
  \beta_0^{(l)}
  + \beta_1^{(l)}x_{1i}
  + \cdots
  + \beta_p^{(l)}x_{pi}
  + u_{j[i]}^{(l)}
  + e_i^{(l)}
  + m_i,
  \qquad \text{(location sub-model)}
  \label{eq:box3_location}\\
  u_j^{(l)}
  &\sim
  \mathcal{N}(0,\sigma_{u,j}^{2}),
  \label{eq:box3_study}\\
  e_i^{(l)}
  &\sim
  \mathcal{N}(0,\sigma_{e,i}^{2}),
  \label{eq:box3_residual}\\
  m_i
  &\sim
  \mathcal{N}(0,v_i),
  \label{eq:box3_sampling}
\end{align}
where $\sigma_{u,j}$ represents between-study heterogeneity,
$\sigma_{e,i}$ represents residual within-study heterogeneity, and $v_i$ is the
known sampling variance. As in the model above, the sampling term can be
generalised to a known variance--covariance matrix when sampling errors are
correlated. Each source of heterogeneity can then be modelled with its own scale
sub-model:
\begin{align}
  \ln(\sigma_{u,j})
  &=
  \beta_0^{(s_u)}
  + \beta_1^{(s_u)}z_{1j}
  + \cdots
  + \beta_q^{(s_u)}z_{qj},
  \qquad \text{(between-study scale sub-model)}
  \label{eq:box3_scale1}\\
  \ln(\sigma_{e,i})
  &=
  \beta_0^{(s_e)}
  + \beta_1^{(s_e)}w_{1i}
  + \cdots
  + \beta_r^{(s_e)}w_{ri},
  \qquad \text{(within-study scale sub-model)}
  \label{eq:box3_scale2}
\end{align}
where the moderators in the three sub-models need not be the same. The location
coefficients, $\beta^{(l)}$, describe how moderators shift the expected effect
size; $\beta^{(s_u)}$ describes how moderators change heterogeneity among
studies; and $\beta^{(s_e)}$ describes how moderators change residual
heterogeneity among effect sizes within studies.

The scale coefficients are additive effects on the natural logarithm of an SD.
For a binary plant-versus-animal moderator, with animal-based fertiliser as the
reference category, the corresponding contrasts on the natural SD scale are
\[
  R_u=\frac{\sigma_{u,P}}{\sigma_{u,A}}
      =\exp\!\left(\beta_{\texttt{plant-animal}}^{(s_u)}\right),
  \qquad
  R_e=\frac{\sigma_{e,P}}{\sigma_{e,A}}
      =\exp\!\left(\beta_{\texttt{plant-animal}}^{(s_e)}\right).
\]
We report these SD ratios because they directly express the multiplicative
difference in heterogeneity: a ratio of 1 indicates equal SDs, and a ratio above
1 indicates greater heterogeneity under plant-based fertiliser. When software
models $\ln(\sigma^2)$ rather than $\ln(\sigma)$, the corresponding SD ratio is
$\exp(\beta/2)$.

This extension provides a more refined view of generalisability and
transferability. Two contexts may have similar expected effects but differ in
where heterogeneity arises. Greater between-study heterogeneity means that
study-level effects vary more strongly around the expected effect, making that
expectation less transferable to a new study. Greater within-study
heterogeneity instead means that individual effect sizes are less consistent
within studies, even if study-level averages remain similarly transferable. A
moderator can therefore affect the location component, either scale component,
or all three simultaneously.

Our organic--conventional yield example illustrates these different
possibilities. For lnRR, between-study heterogeneity was greater under plant-
than animal-based fertilisers (plant-to-animal SD ratio $=2.18$, 95\% CrI
$=[1.11, 5.55]$), whereas the difference in residual heterogeneity was
uncertain (SD ratio $=1.14$, 95\% CrI $=[0.94, 1.40]$). For lnCVR, the pattern
was different: residual heterogeneity was greater under plant-based fertilisers
(SD ratio $=1.67$, 95\% CrI $=[1.18, 2.34]$), whereas the difference in
between-study heterogeneity was highly uncertain (SD ratio $=1.05$, 95\% CrI
$=[0.10, 8.63]$). Thus, the same moderator can influence different levels of
heterogeneity for different ecological outcomes.

Importantly, the extension preserved the broad finding that context dependence
differed between lnRR and lnCVR, while showing that it arose at different
levels. Estimating moderator effects on between-study heterogeneity nevertheless
requires sufficient replication at the study level. In our example, only six
studies contained both fertiliser categories, limiting information about
category-specific study effects. We therefore treat this extension primarily
as an illustration of how multilevel location--scale models can separate
different sources of heterogeneity. Full model specifications, assumptions,
sensitivity analyses, and prediction intervals for new studies are provided in
the online tutorial.

\end{tcolorbox}
